
\documentclass{article}
\usepackage{colortbl}
\usepackage{makecell}
\usepackage{multirow}
\usepackage{supertabular}

\begin{document}

\newcounter{utterance}

\centering \large Interaction Transcript for game `wordle\_withcritic', experiment `medium\_frequency\_words\_clue\_with\_critic', episode 10 with reflexion{-}Qwen3.5{-}9B.
\vspace{24pt}

{ \footnotesize  \setcounter{utterance}{1}
\setlength{\tabcolsep}{0pt}
\begin{supertabular}{c@{$\;$}|p{.15\linewidth}@{}p{.15\linewidth}p{.15\linewidth}p{.15\linewidth}p{.15\linewidth}p{.15\linewidth}}
    \# & $\;$A & \multicolumn{4}{c}{Game Master} & $\;\:$B\\
    \hline

    \theutterance \stepcounter{utterance}  
    & & \multicolumn{4}{p{0.6\linewidth}}{
        \cellcolor[rgb]{0.9,0.9,0.9}{
            \makecell[{{p{\linewidth}}}]{
                \texttt{\tiny{[P1$\langle$GM]}}
                \texttt{You have attempted this same task before and did not succeed. These are the plans you wrote after those attempts.} \\
\texttt{Trial \#1: New plan: I failed because I guessed the clue word "tithing" directly, which is a valid move but likely not the target word, and more critically, I did not explicitly acknowledge the specific rules regarding the Game Master's feedback mechanism in my explanation. The rules state the GM will indicate agreement/disagreement and provide rationale; I must wait for this specific feedback before making my next move. Next time, I will start with a high{-}frequency letter guess (like "crane" or "slate") to gather data, and in my explanation, I will explicitly state that I am waiting for the GM's agreement/disagreement and rationale before proceeding, ensuring I do not generate any unauthorized feedback or skip the required interaction step.} \\
\\ 
\\ 
\texttt{You are a language wizard who likes to guess words by using the given rules.} \\
\\ 
\texttt{Welcome to Wordle! You have six attempts to guess the target word, a valid English word of five lowercase letters (a{-}z). Please use the tags "explanation:" and "guess:" to provide a concise explanation for each guess.} \\
\\ 
\texttt{To help you make an informed guess, you will receive a clue for the word, such as} \\
\texttt{clue: "snowy white"} \\
\\ 
\texttt{Here is an example guess based on the clue:} \\
\texttt{explanation: In the fairy tail Snow White, the girl is killed because she eats a poisoned apple. And the word apple has 5 letters.} \\
\texttt{guess: apple} \\
\\ 
\texttt{I will then indicate whether I agree or disagree with your guess and provide rationale, but agreeing with a guess does not confirm its correctness. You may choose to retain your original guess or modify it based on given clue and agreement.} \\
\\ 
\texttt{After each guess, your answer will be validated, and you will receive feedback indicating which letters are correct (green), which letters are correct but in the wrong position (yellow), and which letters are incorrect (red). This feedback can be useful in determining which letters to include or exclude in your next guess.} \\
\\ 
\texttt{For example, the feedback for "apple" might be:} \\
\texttt{guess\_feedback: a<yellow> p<yellow> p<green> l<yellow> e<red>} \\
\\ 
\texttt{Here are some rules to keep in mind:} \\
\texttt{1. Please do not generate any guess feedback or clue on your own} \\
\texttt{2. Use the clue, my agreement or disagreement, and guess feedback to arrive at your next guess. This should be reflected in your explanation as well.} \\
\\ 
\texttt{Do not generate any other text.} \\
\\ 
\texttt{Let's begin with your first guess.} \\
\\ 
\texttt{clue: tithing} \\
            }
        }
    }
    & \\ \\

    \theutterance \stepcounter{utterance}  
    & \multicolumn{4}{p{0.6\linewidth}}{
        \cellcolor[rgb]{0.9,0.9,0.9}{
            \makecell[{{p{\linewidth}}}]{
                \texttt{\tiny{[P1$\rangle$GM]}}
                \texttt{explanation: The clue "tithing" suggests a word related to giving a tenth part, often associated with religious or charitable giving. A common five{-}letter word fitting this theme is "tithe".} \\
\texttt{guess: tithe} \\
            }
        }
    }
    & & \\ \\

    \theutterance \stepcounter{utterance}  
    & & \multicolumn{4}{p{0.6\linewidth}}{
        \cellcolor[rgb]{0.9,0.9,0.9}{
            \makecell[{{p{\linewidth}}}]{
                \texttt{\tiny{[GM$\rangle$P2]}}
                \texttt{I need your assistance with a word game in which we need to find a 5{-}letter word using a clue, a guess and an explanation for the guess.} \\
\\ 
\texttt{For example, if the clue is "snow white" and the guess is "apple", you will get the information like this:} \\
\texttt{clue: snow white} \\
\texttt{explanation: the fairy tale Snow White involves the girl eating a poisoned apple, which has 5 letters.} \\
\texttt{guess: apple} \\
\\ 
\texttt{Your task is to either agree or disagree with my guess based on the given clue.} \\
\\ 
\texttt{If you disagree, I will make another guess, and if you agree, I will proceed in the game using the same guess.} \\
\\ 
\texttt{Initially, guess feedback is not available, but as the game proceeds, it becomes available. The guess feedback is not for current guess, but rather an indication of what letters can be present in the current guess. A letter highlighted in green indicates that it is a correct letter in the correct position, while yellow indicates that it is a correct letter in the wrong position, and red indicates a wrong letter. At this point, you should use all the provided information, including the clue, guess, explanation, and guess feedback, to decide whether the given guess matches the clue, and also stick to guess feedback or not.} \\
\\ 
\texttt{Please respond in lowercase letters and stick to this format:} \\
\texttt{explanation: your reason for doing so} \\
\texttt{agreement: yes or no} \\
\\ 
\texttt{Please note that you must not guess any word based on the riddle clue and stick to the given format while responding.} \\
\\ 
\texttt{Do not generate any other text.} \\
\\ 
\texttt{Let's begin!} \\
\\ 
\texttt{clue: tithing} \\
\texttt{explanation: The clue "tithing" suggests a word related to giving a tenth part, often associated with religious or charitable giving. A common five{-}letter word fitting this theme is "tithe".} \\
\texttt{guess: tithe} \\
            }
        }
    }
    & \\ \\

    \theutterance \stepcounter{utterance}  
    & & & \multicolumn{4}{p{0.6\linewidth}}{
        \cellcolor[rgb]{0.9,0.9,0.9}{
            \makecell[{{p{\linewidth}}}]{
                \texttt{\tiny{[GM$\langle$P2]}}
                \texttt{explanation: the word "tithe" directly matches the clue "tithing" as it is the specific term for giving a tenth part, and it has exactly 5 letters.} \\
\texttt{agreement: yes} \\
            }
        }
    }
     \\ \\

    \theutterance \stepcounter{utterance}  
    & & \multicolumn{4}{p{0.6\linewidth}}{
        \cellcolor[rgb]{0.9,0.9,0.9}{
            \makecell[{{p{\linewidth}}}]{
                \texttt{\tiny{[P1$\langle$GM]}}
                \texttt{clue: tithing} \\
\texttt{guess\_agreement: yes} \\
\texttt{agreement\_explanation: the word "tithe" directly matches the clue "tithing" as it is the specific term for giving a tenth part, and it has exactly 5 letters.} \\
\\ 
\texttt{Provide your response only in this format.} \\
\texttt{explanation: details} \\
\texttt{guess: word} \\
            }
        }
    }
    & \\ \\

    \theutterance \stepcounter{utterance}  
    & \multicolumn{4}{p{0.6\linewidth}}{
        \cellcolor[rgb]{0.9,0.9,0.9}{
            \makecell[{{p{\linewidth}}}]{
                \texttt{\tiny{[P1$\rangle$GM]}}
                \texttt{explanation: The GM agreed that "tithe" matches the clue "tithing" and has 5 letters, but agreement does not confirm correctness. I must wait for the specific guess feedback (green/yellow/red) before making my next move to validate the word.} \\
\texttt{guess: tithe} \\
            }
        }
    }
    & & \\ \\

    \theutterance \stepcounter{utterance}  
    & & & \multicolumn{2}{p{0.3\linewidth}}{
        \cellcolor[rgb]{0.9,0.9,0.9}{
            \makecell[{{p{\linewidth}}}]{
                \texttt{\tiny{[GM$|$GM]}}
                \texttt{attempts = 1} \\
\texttt{target\_word = tenth} \\
\texttt{target\_word\_clue = tithing} \\
\texttt{guess = tithe} \\
\texttt{guess\_feedback = t<green> i<red> t<yellow> h<yellow> e<yellow>} \\
            }
        }
    }
    & & \\ \\

    \theutterance \stepcounter{utterance}  
    & & \multicolumn{4}{p{0.6\linewidth}}{
        \cellcolor[rgb]{0.9,0.9,0.9}{
            \makecell[{{p{\linewidth}}}]{
                \texttt{\tiny{[P1$\langle$GM]}}
                \texttt{guess\_feedback: t<green> i<red> t<yellow> h<yellow> e<yellow>} \\
\\ 
\texttt{Provide your response only in this format.} \\
\texttt{explanation: details} \\
\texttt{guess: word} \\
            }
        }
    }
    & \\ \\

    \theutterance \stepcounter{utterance}  
    & \multicolumn{4}{p{0.6\linewidth}}{
        \cellcolor[rgb]{0.9,0.9,0.9}{
            \makecell[{{p{\linewidth}}}]{
                \texttt{\tiny{[P1$\rangle$GM]}}
                \texttt{explanation: The feedback shows 't' is correct in the first position (green), 'i' is incorrect (red), and 't', 'h', 'e' are correct but in wrong positions (yellow). Since 't' is already in the first position, the second 't' must be in the third or fourth position. The letters 'h' and 'e' need to be moved. A strong candidate fitting this pattern is "theft" (t{-}h{-}e{-}f{-}t), but let's try "thief" (t{-}h{-}i{-}e{-}f) or "theft". Wait, 'i' is red, so 'thief' is out. How about "thems"? No. Let's try "theft": t (green), h (yellow), e (yellow), f (unknown), t (yellow). If the word is "theft", the second 't' would be green if it's the last letter, but here the last letter is yellow. So the second 't' is not at the end. Maybe "thigh"? t (green), h (yellow), i (red {-} conflict), g (unknown), h (yellow). 'i' is red, so no 'i'. How about "thine"? t (green), h (yellow), i (red {-} conflict). How about "thems"? t (green), h (yellow), e (yellow), m (unknown), s (unknown). This fits if 't' is only at start and 'e' is in middle. But we have two 't's in the feedback (first green, third yellow). So the word has two 't's. One at start, one elsewhere. 'i' is out. 'h' and 'e' are misplaced. Let's try "thete"? No. "thete" isn't a word. "thete" {-}> t{-}h{-}e{-}t{-}e. t (green), h (yellow), e (yellow), t (yellow), e (yellow). This fits perfectly! t is at 1 and 4. h is at 2 (wrong pos). e is at 3 and 5 (wrong pos). Wait, if the word is "thete", then 'e' is at 3 and 5. In my guess "tithe", 'e' was at 5 (yellow). In "thete", 'e' would be at 3 and 5. But I only have one 'e' in my guess. The feedback says 'e' is yellow, meaning it's in the word but wrong spot. If the word is "thete", 'e' is at 3 and 5. My guess had 'e' at 5. So 'e' at 5 is yellow (correct letter, wrong spot). 'e' at 3 is also yellow? No, I only guessed one 'e} \\
            }
        }
    }
    & & \\ \\

    \theutterance \stepcounter{utterance}  
    & & & \multicolumn{2}{p{0.3\linewidth}}{
        \cellcolor[rgb]{0.9,0.9,0.9}{
            \makecell[{{p{\linewidth}}}]{
                \texttt{\tiny{[GM$|$GM]}}
                \texttt{Error: The response should contain the 'guess:' keyword exactly once.} \\
            }
        }
    }
    & & \\ \\

    \theutterance \stepcounter{utterance}  
    & & & \multicolumn{2}{p{0.3\linewidth}}{
        \cellcolor[rgb]{0.9,0.9,0.9}{
            \makecell[{{p{\linewidth}}}]{
                \texttt{\tiny{[GM$|$GM]}}
                \texttt{game\_result = ABORT} \\
            }
        }
    }
    & & \\ \\

    \theutterance \stepcounter{utterance}  
    & \multicolumn{4}{p{0.6\linewidth}}{
        \cellcolor[rgb]{0.95,0.95,0.95}{
            \makecell[{{p{\linewidth}}}]{
                \texttt{\tiny{[P2$\rightarrow$P1]}}
                \texttt{Post{-}episode Reflexion (not sent to GM):} \\
\texttt{New plan: I failed because my explanation became a rambling internal monologue that violated the "Do not generate any other text" rule and did not strictly adhere to the required format. Specifically, I listed multiple invalid word candidates ("thete", "thine") and failed to provide a single, definitive guess within the `guess:` tag. Next time, I will analyze the feedback immediately to deduce the correct word structure (e.g., T \_ \_ \_ E with specific letter placements), select exactly one valid English word that fits the constraints, and output a concise explanation followed by that single word in the correct format without any extra commentary or invalid options.} \\
            }
        }
    }
    & & \\ \\

    \theutterance \stepcounter{utterance}  
    & & & \multicolumn{4}{p{0.6\linewidth}}{
        \cellcolor[rgb]{0.95,0.95,0.95}{
            \makecell[{{p{\linewidth}}}]{
                \texttt{\tiny{[P1$\rightarrow$P2]}}
                \texttt{Post{-}episode Reflexion (not sent to GM):} \\
\texttt{New plan: I failed because I violated the explicit rule "must not guess any word based on the riddle clue" by providing an explanation that confirmed the guess "tithe" was correct. The task requires me to act as a validator who either agrees or disagrees with the user's guess based on the clue, not to generate the answer myself. Next time, I will strictly evaluate the user's provided guess against the clue without revealing the solution or confirming the guess is correct in my explanation. I will only output "agreement: yes" if the user's guess perfectly matches the clue's meaning and length, or "agreement: no" if it does not, ensuring my explanation describes the validation logic rather than solving the riddle.} \\
            }
        }
    }
     \\ \\

\end{supertabular}
}

\end{document}
